\documentclass[addpoints]{exam}
% \documentclass[addpoints,12pt]{exam}

\usepackage[slovene]{babel}
\usepackage{amsfonts}
\usepackage{amssymb}
\usepackage{amsmath}
\usepackage{float}
\usepackage{graphicx}
\usepackage{enumitem}
\usepackage{color}
\usepackage{eurosym}


%Komentar naslednje vrstice glede na verzijo testa -- rešitve ali prostor za odgovore:
\printanswers

% \linespread{2}


\pointsinrightmargin
\marginbonuspointname{}

\renewcommand{\solutiontitle}{\noindent\textbf{Rešitve in točkovanje:}\par\noindent}

\colorgrids
\definecolor{GridColor}{gray}{0.7}
\setlength{\gridsize}{7mm}

\extrawidth{5mm}
\setlength{\rightpointsmargin}{1.5cm}

\begin{document}

\runningfooter{}{}{}

\begin{center}
    \fbox{\fbox{\parbox{15cm}{\centering
    Prvo pisno ocenjevanja znanja \\ 1. letnik -- test B \\ 23. oktober 2024 \\ (Osnove logike in teorije množic, Naravna števila)}}}
\end{center}

\vspace{0.1in}
\makebox[\textwidth]{Ime in priimek, oddelek:\enspace\hrulefill}


\begin{center}
    \chqword{Naloga}
    \chpword{Možne točke}
    \chbpword{Dodatne točke}
    \chsword{Dosežene točke}
    \chtword{Skupaj}  
    \combinedpointtable[h][questions]
    % \combinedgradetable[h][questions]
\end{center}

\vspace{0.1in}


\begin{questions}

    % 1. vprašanje
    
    \question[4]
    Določite pravilnost izjave $(\lnot A\Rightarrow  B)\lor (A\Leftrightarrow \lnot B)$ glede na izbor izjav $A$ in $B$.

    \begin{solution}[8cm]
        \begin{table}[H]
            \begin{tabular}{ccccccc}
             $A$ & $B$ & $\lnot A$ & $\lnot B$ & $(\lnot A\Rightarrow B)$ & $(A\Leftrightarrow \lnot B)$ & $(\lnot A\Rightarrow B)\lor(A\Leftrightarrow \lnot B)$\\ \hline
             P & P & N & N & P & N & P \\ 
             P & N & N & P & P & P & P \\
             N & P & P & N & P & P & P \\
             N & N & P & P & N & N & N
            \end{tabular}
        \end{table}
        Za vsako pravilno izpolnjeno vrstico, ki določa pravilnost, po $1$ točka.

    \end{solution}




    % 2. vprašanje
    \question[3]
    Določite resničnostno vrednost izjav:
    \begin{parts}
        \part 
        Število $23$ ni večkratnik števila 13 ali število $9$ je liho praštevilo.\\
        \part 
        Če število $3$ ni praštevilo, potem je praštevil neskončno mnogo.\\
        \part
        Ni res, da je ostanek pri deljenju s $4$ lahko samo $1$, $2$ ali $3$ in število $8$ je večje od števila $2$.\\
    \end{parts}

    \begin{solution}
        (a): P (prva izjava je pravilna, druga je nepravilna) \\
        (b): P (prva izjava je nepravilna, druga je pravilna) \\
        (c): P (negirana izjava je pravilna /prva izjava je nepravilna, druga izjava je pravilna/) \\ \\
        Za vsako pravilno določeno logično vrednost po $1$ točka.
    \end{solution}



    \newpage

    % 3. vprašanje
    \question
    Dani sta množici $\mathcal{A}=\left\{e, f, g, m\right\}$ in $\mathcal{B}=\left\{f, g, m, u\right\}$.
    \begin{parts}
        \part[2]
        Ali velja $\mathcal{A}\subseteq\mathcal{B}$? Poiščite največjo možno množico $\mathcal{C}$, da bosta množici $\mathcal{A}$ in $\mathcal{B}$ njeni nadmnožici.
        \part[3]
        Zapišite $\mathcal{PA}$. Izračunajte koliko elementov vsebuje.
        \part[2]
        Zapišite in grafično predstavite kartezični produkt $\mathcal{A}\times\mathcal{B}$.
    \end{parts}

    % \begin{description}
    %     \item[a] $(x-a)^2+(y-b)^2-c^2=0$ 
    %     \item[b] $d^2(x-a)^2+f^2(y-b)^2-(df)^2=0$
    %     \item[c] $S(a,b)$
    %     \item[d] $e^2=d^2-f^2;\ d>f$
    % \end{description}
    % Krožnica: \fillin[a c][2cm]
    % \\ Elipsa:  \fillin[b c d][2cm]

    \begin{solution}[9cm]
        (a): Ne drži, da $\mathcal{A}\subseteq\mathcal{B}$, ker $e\in\mathcal{A}\land e\notin\mathcal{B}$.; $\mathcal{C}=\left\{f,g,m\right\}$ \\
        (b): $\mathcal{PA}=\left\{\emptyset, \{e\}, \{f\}, \{g\}, \{m\}, \{e,f\}, \{e,g\}, \{e,m\}, \{f,g\}, \{f,m\}, \{g,m\}, \{e,f,g\}, \{e,f,m\}, \{e,g,m\}, \{f,g,m\}, \mathcal{A}\right\}$; $m(\mathcal{PA})=2^{m(\mathcal{A})}=2^4=16$. \\
        (c): $\mathcal{A}\times\mathcal{B}=\left\{(e,f), (e,g), (e,m), (e,u), (f,f), (f,g), (f,m), (f,u), (g,f), (g,g), (g,m), (g,u), (m,f), (m,g), (m,m), (m,u)\right\}$ + grafična upodobitev \\ \\
        Pri (a) za vsak odgovor po $1$ točka. Pri (b) za zapis celotne $\mathcal{PA}$ $2$ točki, za zapis zgolj podmnožic $1$ točka; za izračun po formuli za moč potenčne množice $1$ točka. Pri (c) za zapis $\mathcal{A}\times\mathcal{B}$ $1$ točka, za grafično upodobitev $1$ točka.
    \end{solution}


    % 4. vprašanje
    \question
    Množica $\mathcal{C}$ je prava podmnožica množice $\mathcal{B}$, prav tako je $\mathcal{B}$ prava podmnožica množice $\mathcal{A}$.
    \begin{parts}
        \part[1]
        Grafično predstavite množico $(\mathcal{A}\setminus\mathcal{B})\cup\mathcal{C}$.
        \part[3]
        Izračunajte $\mathcal{B}\setminus\mathcal{A}$, $\mathcal{C}\cap\mathcal{A}$ in $\mathcal{C}\cup\mathcal{B}$.
    \end{parts}

    \begin{solution}[1cm]
        (a): grafična upodobitev \\
        (b): $\mathcal{B}\setminus\mathcal{A}=\emptyset$; $\mathcal{C}\cap\mathcal{A}=\mathcal{C}$; $\mathcal{C}\cup\mathcal{B}=\mathcal{B}$ \\ \\
        Pri (a) za narisan Euler-Vennov diagram $1$ točka. Pri (b) za izračun vsake izmed množic po $1$ točka.
    \end{solution}




    % \newpage

    % 5. vprašanje
    \question
    Naj bo $\mathcal{U}=\mathbb{N}_{22}$. Dane so množice $\mathcal{A}=\left\{x;\ x=2n \land 3\leq n<10\right\}$, $\mathcal{B}=\left\{x;\ x=3k-2 \land 3\leq k\leq 7\right\}$ in $\mathcal{C}=\left\{x;\ x=5m \land 1<m\leq 4\right\}$. 
    \begin{parts}
        \part[5]
        Določite naslednje množice: $\mathcal{A}\cap\mathcal{C}$, $\mathcal{C}\setminus\mathcal{A}$, $(\mathcal{A}\cup\mathcal{C})\setminus\mathcal{B}$ in $(\mathcal{A}\cup\mathcal{B})^\complement$.
        \part[1]
        Koliko elementov ima množica $\mathcal{C}\times\mathcal{B}$?
    \end{parts}

    \begin{solution}[9cm]
        $\mathcal{A}=\left\{6,8,10,12,14,16,18\right\}$, $\mathcal{B}=\left\{7,10,13,16,19\right\}$, $\mathcal{C}=\left\{10,15,20\right\}$ \\
        (a): $\mathcal{A}\cap\mathcal{C}=\left\{10\right\}$; $\mathcal{C}\setminus\mathcal{A}=\left\{15,20\right\}$; $(\mathcal{A}\cup\mathcal{C})\setminus\mathcal{B}=\left\{6,8,12,14,15,18,20\right\}$; $(\mathcal{A}\cup\mathcal{B})^\complement=\left\{1,2,3,4,5,9,11,15,17,20,21,22\right\}$ \\
        (b): $m(\mathcal{C}\times\mathcal{B})=m(\mathcal{C})\cdot m(\mathcal{B})=3\cdot 5 = 15$ \\ \\
        Pri (a) za zapis/uporabo pravilnih množic $\mathcal{A}, \mathcal{B}, \mathcal{C}$ $1$ točka; za vsako pravilno določeno množico po $1$ točka. Pri (b) za uporabo formule in izračun moči kartezičnega produkta $1$ točka.
    \end{solution}


\newpage
    % 6. vprašanje
    \question
    Pri uri športa so se dijaki in učitelj pogovarjali o žogah.  Učitelj je dijake vprašal, katere vrste žog imajo doma.
    $20$ dijakov ima \textit{nogometno žogo}, $6$ jih ima \textit{košarkarsko žogo} in $17$ \textit{rokometno žogo}. Kar $12$ jih ima \textit{nogometno} in \textit{rokometno žogo}, $4$ imajo \textit{košarkarsko} in \textit{nogometno žogo},
    $3$ pa \textit{košarkarsko} in \textit{rokometno žogo}. En sam dijaki ima vse tri vrste žog, dva pa nimata doma nobene žoge. 
    \begin{parts}
        \part[1]
        Narišite diagram.
        \part[2]
        Koliko dijakov je bilo pri uri športa?
        \part[2]
        Koliko dijakov ima doma natanko eno žogo?
    \end{parts}

    \begin{solution}[8cm]
        (a): Euler-Vennov diagram z vpisanimi številkami \\
        (b): $5+0+3+3+2+11+1+2=27$ dijakov \\
        (c): $5+3+0=8$ dijakov \\ \\
        Pri (a) za pravilno narisan in izpolnjen diagram $1$ točka. Pri (b) in (c) za pravilen izračun $1$ točka, za zapis odgovora $1$ točka.
    \end{solution}




    % \newpage
        

    % 7. vprašanje
    \question[3]
    Na šoli je šest oddelkov s po $24$ dijaki, pet oddelkov s po $32$ dijaki in štirje oddelki s po $29$ dijaki. 
    Med vsemi dijaki je $315$ fantov. Koliko deklet je na šoli? (Zapišite samo en račun.)

    \begin{solution}[7cm]
        $6\cdot 24+5\cdot 32+4\cdot 29-315=105$ \\ \\
        Za nastavek enega računa $1$ točka, za pravilen izračun $1$ točka, za zapis odgovora $1$ točka.
    \end{solution}

    
    %dodatno vprašanje
    \bonusquestion[3]
    \textit{Dodatna naloga:} Izjavo \textit{Obstaja takšno naravno število $x$, da deli število $5$ ali število $7$.} zapišite s simboli in jo zanikajte v simbolnem zapisu.
    
    \begin{solution}
        $\exists x\in\mathbb{N}\ni:(x\mid 5)\lor(x\mid 7)$; negacija: $\forall x\in\mathbb{N}: (x\nmid 5)\land(x\nmid 7)$ \\ \\
        Za pravilen zapis izjave s simboli $1$ točka, za pravilno negacijo kvantifikatorja $1$ točka, za pravilno negacijo logične operacije $1$ točka.
    \end{solution}

\end{questions}



\end{document}